%% =====================================================================
%%  Corrigés — Unit 3 « What's your number? »
%% =====================================================================
\corrigesUnite{3}{What's your number?}

\begin{corrige}{1}
\begin{reponses}[4]
\item \rep{seven}
\item \rep{twelve}
\item \rep{fifteen}
\item \rep{twenty}
\item \rep{thirty-three}
\item \rep{forty-eight}
\item \rep{fifty}
\item \rep{ninety-nine}
\end{reponses}
\piege{écrire \en{fourty} : le \en{u} de \en{four} disparaît dans \en{forty}. Et n'oubliez pas le
trait d'union de \en{thirty-three}, \en{forty-eight}, \en{ninety-nine}.}
\end{corrige}

\begin{corrige}{2}
\begin{reponses}[6]
\item \rep{13}
\item \rep{30}
\item \rep{81}
\item \rep{66}
\item \rep{14}
\item \rep{40}
\end{reponses}
\retiens{\en{-teen} = de 13 à 19 (\en{thirteen}, \en{fourteen}) ; \en{-ty} = les dizaines
(\en{thirty}, \en{forty}). À l'oral, l'accent aide : \en{thir\textbf{teen}}, mais \en{\textbf{thir}ty}.}
\end{corrige}

\begin{corrige}{3}
\begin{reponses}[4]
\item \rep{a}
\item \rep{an}
\item \rep{a}
\item \rep{an}
\item \rep{a}
\item \rep{an}
\item \rep{a}
\item \rep{an}
\end{reponses}
\pourquoi{c'est le \textbf{son} qui compte : \en{hour} commence par un son voyelle (le \en{h} est
muet), \en{university} par le son \phon{j} de \en{you}. Au numéro 3, c'est \en{big} qui suit
\en{a} ; au numéro 8, c'est \en{old}, qui commence par une voyelle.}
\end{corrige}

\begin{corrige}{4}
\begin{reponses}[3]
\item \rep{What}
\item \rep{Where}
\item \rep{How old}
\item \rep{How}
\item \rep{Who}
\item \rep{What}
\end{reponses}
\pourquoi{la réponse guide le choix : un lieu → \en{Where} ; un âge → \en{How old} ; une
personne → \en{Who} ; un état (\en{fine}) → \en{How} ; un nom ou un numéro → \en{What}.}
\end{corrige}

\begin{corrige}{5}
\begin{reponses}[2]
\item \rep{How old are you?}
\item \rep{How do you spell it?}
\item \rep{What's your phone number?}
\item \rep{I'm twenty-seven.}
\item \rep{Who's that?}
\end{reponses}
\end{corrige}

\begin{corrige}{6}
\begin{reponses}[2]
\item \rep{It's Brooks.}
\item \rep{I'm twenty-four.}
\item \rep{It's 718-555-0162.}
\item \rep{No, I'm not. I'm American.}
\end{reponses}
\pourquoi{pour un nom ou un numéro, on répond avec \en{It's…} ; pour l'âge, avec \en{I'm…}. À la
question 4, la réponse courte est suivie de la bonne information.}
\end{corrige}

\begin{corrige}{7}
\begin{reponses}[3]
\item \rep{Carter}
\item \rep{Lopez}
\item \rep{Brooks}
\item \rep{Hughes}
\item \rep{Gray}
\item \rep{Quinn}
\end{reponses}
\piege{confondre \en{E} \phon{iː} et \en{I} \phon{aɪ}, ou \en{G} \phon{dʒiː} et \en{J}
\phon{dʒeɪ}. Dans \en{Hughes}, on entend \en{H, U, G, H, E, S}. Texte de la piste
page~\pageref{trans-3.11}.}
\end{corrige}

\begin{corrige}{8}
\begin{reponses}[2]
\item \rep{212-555-0134}
\item \rep{31}
\item \rep{347-555-0119}
\item \rep{45}
\item \rep{646-555-0172}
\item \rep{23}
\item \rep{718-555-0156}
\item \rep{38}
\end{reponses}
\retiens{dans un numéro, \en{oh} = 0. Distinguez bien \en{thirty-one} (31) et \en{thirteen} (13).}
\end{corrige}

\begin{corrige}{9}
\begin{reponses}[1]
\item \rep{jack.miller@example.com}
\item \rep{emma\_clarke@example.co.uk}
\item \rep{studio@example.org}
\item \rep{sam.brooks24@example.com}
\end{reponses}
\pourquoi{\en{at} = @, \en{dot} = point, \en{underscore} = tiret bas. \en{co.uk} est la terminaison
des adresses britanniques.}
\end{corrige}

\begin{corrige}{10}
\begin{reponses}[1]
\item \rep{I'm twenty-seven.} \textit{ou} \rep{I'm twenty-seven years old.}
\item \rep{How do you spell it?}
\item \rep{What's your phone number?}
\item \rep{How old is Jack? — He's thirty.}
\item \rep{It's an email address.}
\item \rep{Who's that? — That's Emma.}
\item \rep{Sorry? One more time, please.}
\item \rep{My last name is Martin.}
\end{reponses}
\piege{traduire « j'ai … ans » avec \en{have}. L'âge s'exprime toujours avec \en{be}.}
\end{corrige}

\begin{corrige}{11}
Réponses modèles, avec les informations de Léa :
\begin{reponses}[1]
\item \rep{My first name is Léa.}
\item \rep{My last name is Martin. M, A, R, T, I, N.}
\item \rep{I'm twenty-seven.}
\item \rep{I'm from Lyon, in France.}
\item \rep{It's 917-555-0148.}
\item \rep{It's lea.martin@example.com.}
\end{reponses}
\retiens{vérifiez trois points : la majuscule des noms propres, \en{I'm} pour l'âge, \en{It's}
pour un numéro ou une adresse.}
\end{corrige}

\begin{corrige}{12}
Réponses modèles :
\begin{reponses}[1]
\item \rep{My name is Léa Martin.}
\item \rep{M, A, R, T, I, N.}
\item \rep{I'm twenty-seven.}
\item \rep{It's 917-555-0148.}
\item \rep{Thank you! You too!} \textit{ou} \rep{Thanks! Have a nice day!}
\end{reponses}
\end{corrige}

\begin{corrigetest}
\begin{reponses}[5][french]
\item \rep{b}
\item \rep{c}
\item \rep{a}
\item \rep{b}
\item \rep{c}
\item \rep{a}
\item \rep{b}
\item \rep{a}
\item \rep{a}
\item \rep{c}
\end{reponses}
\pourquoi{question 1 : l'âge se dit avec \en{be}. Questions 4 et 5 : \en{an} devant le son voyelle
de \en{email}, \en{a} devant le son consonne de \en{phone}.}
\end{corrigetest}
